% Generated draft: requires mathematical and source review.
\documentclass[11pt]{article}
\usepackage[margin=25mm]{geometry}
\usepackage{fontspec}
\setmainfont{Latin Modern Roman}
\usepackage{amsmath,amssymb,mathtools,unicode-math}
\setmathfont{Latin Modern Math}
\usepackage{enumitem,needspace,xurl,hyperref}
\setlength{\parindent}{0pt}
\setlength{\parskip}{5pt}
\setlength{\emergencystretch}{3em}
\newcommand{\N}{\mathbb{N}}
\newcommand{\Z}{\mathbb{Z}}
\newcommand{\Q}{\mathbb{Q}}
\newcommand{\R}{\mathbb{R}}
\newcommand{\C}{\mathbb{C}}
\newcommand{\F}{\mathbb{F}}
\newcommand{\A}{\mathbb{A}}
\newcommand{\PP}{\mathbb{P}}
\newcommand{\E}{\mathbb{E}}
\newcommand{\eps}{\varepsilon}
\newcommand{\dd}{\,\mathrm{d}}
\begin{document}
\section*{UnsolvedMath ID: 700904}
\textbf{Source problem code:} catalogue entry
\subsection{Definitions and mathematical statement}
Let \(\Gamma=(V,E)\) be a nonempty connected undirected simple graph: \(E\) consists of unordered pairs of distinct vertices, with no multiple edges. No finiteness or local finiteness hypothesis is imposed in the supplied statement. Let \(G=\operatorname{Aut}(\Gamma)\), and assume that \(G\) acts transitively on \(V\). For \(v\in V\), write \(G_v=\{g\in G:g(v)=v\}\).

A \(G\)-invariant partition \(\sigma\) of \(V\) is a partition into nonempty blocks such that \(g(B)\in\sigma\) for every \(g\in G\) and \(B\in\sigma\). Its blocks are finite if \(|B|<\infty\) for every \(B\in\sigma\). Define the simple quotient graph \(\Gamma/\sigma\) to have vertex set \(\sigma\), with distinct blocks \(B,C\) adjacent precisely when some vertex of \(B\) is adjacent in \(\Gamma\) to some vertex of \(C\). Edges within a block produce no quotient loops. The source calls \(\sigma\) an imprimitivity system but does not specify whether trivial invariant partitions are excluded.

Let \(d_\Gamma\) denote graph distance. For a positive integer \(n\), let \(\Gamma^{[n]}\) have vertex set \(V\) and edge set
\[
E_n=\bigl\{\{x,y\}:x,y\in V,\ x\ne y,\ d_\Gamma(x,y)\le n\bigr\}.
\]
A regular tree of valency \(3\) means a connected acyclic graph in which every vertex has degree \(3\); such a tree is infinite. Here containment is interpreted as containment as a subgraph, without requiring the subgraph to be induced. The supplied source does not explicitly specify this convention.

The question attributed to V. I. Trofimov is whether every such \(\Gamma\) satisfies at least one of the following assertions:
\begin{enumerate}
\item There exists \(v\in V\) such that \(|G_v|<\infty\).
\item There exists a \(G\)-invariant partition \(\sigma\) into finite blocks and a block \(B\in\sigma\) such that
\[
\bigl|\operatorname{Aut}(\Gamma/\sigma)_B\bigr|<\infty.
\]
The stabilizer here is taken in the full automorphism group of the quotient graph, not merely in the permutation group induced by \(G\).
\item There exists an integer \(n\ge1\) such that \(\Gamma^{[n]}\) contains a regular tree of valency \(3\) as a subgraph.
\end{enumerate}
The conventions concerning imprimitivity systems, quotient graphs, and subgraph containment above make the proposed formulation explicit and require confirmation against the cited source.
\subsection{Short English statement}
Must every connected vertex-transitive simple graph have a finite vertex stabilizer, admit a quotient by finite invariant blocks whose full automorphism group has finite vertex stabilizers, or contain an infinite cubic tree after edges are added between vertices at some uniformly bounded distance? No local finiteness assumption is supplied. The precise conventions for block systems and containment require review.
\subsection{Sources}
\begin{itemize}
\item Local UnsolvedMath registry, permanent ID \(700904\), exact source problem code \texttt{catalogue entry}, title \emph{Kourovka Notebook Problem 12.87}; supplied upstream revision \texttt{d371ee917ac9647fbb06cd6fc529deec6f3dda82}.
\item \emph{The Kourovka Notebook (Unsolved Problems in Group Theory)}, original problem 12.87, 12th issue (1992), proposed by V. I. Trofimov. The supplied record also identifies the 21st issue edition (2026), \texttt{arXiv:1401.0300v48}.
\item The registry describes the collection as Issues 1--21 (1965--2026), edited by E. I. Khukhro and V. D. Mazurov. Its label \texttt{open} and its update date 2026-10-16 are supplied registry metadata, not an independently verified mathematical or literature assessment. No separate literature-verification date or literature-source links were supplied.
\end{itemize}
\end{document}
